\subsection{Zeitsynchronisation über NTP}
Eine zuverlässige Zeitbasis ist für zeitabhängige Systeme erforderlich, da lokale Uhren ohne externe Synchronisation Abweichungen aufweisen können.
Mills beschreibt, dass lokale Uhren häufig auf Quarzoszillatoren basieren, deren Abweichung bis zu etwa einer Sekunde pro Tag betragen kann~\cite{mills1991ntp}.

Das \gls{ntp}, dient zur Verteilung von Zeitinformationen in verteilten Netzwerken~\cite{mills1991ntp}.
Das Protokoll nutzt eine hierarchische Struktur aus primären und sekundären Zeitservern.
Primäre Zeitserver sind direkt mit einer Referenzzeitquelle verbunden.
Sekundäre Zeitserver beziehen ihre Zeit über Netzwerkpfade von primären oder weiteren sekundären Servern.
Die Synchronisation erfolgt über den Austausch von Zeitstempeln zwischen den beteiligten Systemen.
Aus diesen Zeitstempeln werden die Übertragungsverzögerung und die Zeitabweichung zwischen lokaler Uhr und Zeitserver bestimmt~\cite{mills1991ntp}.
Auf dieser Grundlage kann die lokale Uhr korrigiert werden.

Für den Einsatz in Netzwerken berücksichtigt \gls{ntp} schwankende Übertragungsverzögerungen und mögliche Störungen.
Dazu werden Filter-, Auswahl- und Kombinationsverfahren verwendet, um geeignete Zeitinformationen auszuwählen.
Nach Mills kann mit \gls{ntp} in weiten Teilen des Internets eine Zeitgenauigkeit im Bereich weniger Millisekunden erreicht werden~\cite{mills1991ntp}.
Für die entwickelte Wortuhr ist diese Genauigkeit ausreichend, da die Zeitanzeige in diskreten Minutenintervallen erfolgt.
\gls{ntp} stellt damit eine geeignete Grundlage dar, um die interne Zeitbasis des Mikrocontrollers regelmäßig mit einer externen Referenzzeit abzugleichen.
% Network Time Protocol (NTP)
% UTC vs. lokale Zeit
% Zeitzonen und Sommerzeit (DST)
% Zeitserver (pool.ntp.org)
% Synchronisationsstrategien (Start + regelmäßig)
% Fehlerfälle (Timeout, Retry)
% Driftkorrektur